\section{Functions and practical examples}
\subsection{Start with the supplied example}
Call \texttt{App\_Init()} once after HAL, clock and GPIO initialization, then call
\texttt{App\_Process()} repeatedly from the main loop. Set
\texttt{APP\_SENSOR\_COUNT} to one for one sensor or to the required capacity.
\texttt{g\_sensor\_roms[i]}, \texttt{g\_temperature\_c[i]} and
\texttt{g\_temperature\_status[i]} refer to the same slot. Index zero is the first
sensor. Only an OK status permits using that slot's temperature.

For another MCU, bind its GPIO/time callbacks first. The library offers a
complete blocking measurement, separate start/read operations and a cooperative
array interface. The array interface also works with a single sensor.

\subsection{Initialize the bus: OneWire\_Init}
Inputs are the destination bus and a callback configuration. The function
copies a valid configuration, releases the line and returns true; it does not
discover a sensor. The mandatory callbacks and paired critical callbacks were
introduced in the architecture section. The concrete STM32 example is in
\texttt{Application/Src/app.c}.
\begin{lstlisting}
OneWire_Bus bus;
/* callbacks contains your board's GPIO/time functions. */
bool bus_ok = OneWire_Init(&bus, &callbacks);
\end{lstlisting}
Do not use the bus if initialization fails. Keep it alive while the array
refers to it. The following snippets assume initialization succeeded.

\subsection{Obtain an address: ReadRom and IsValidRom}
\texttt{DS18B20\_ReadRom} takes the bus and an eight-byte output buffer and
returns a status. Use it only when exactly one sensor is physically connected.
\texttt{DS18B20\_IsValidRom} checks family and CRC of a supplied address;
it does not prove the sensor is currently present.
\begin{lstlisting}
uint8_t rom[DS18B20_ROM_SIZE];
DS18B20_Status result = DS18B20_ReadRom(&bus, rom);
bool address_ok = (result == DS18B20_STATUS_OK);
\end{lstlisting}
For several sensors, choose the array's automatic discovery or use
\texttt{OneWire\_SearchFirst} and \texttt{OneWire\_SearchNext} with caller-owned
search state. Each successful search supplies a CRC-valid ROM; also check its
DS18B20 family. Search order does not identify a physical location.

\subsection{Read one complete measurement: ReadTemperatureBlocking}
Inputs are bus, ROM, nominal polling budget in milliseconds and output pointer.
The function starts conversion, waits and reads a validated temperature. Only
OK updates the output; on failure it stays unchanged. Use 1000 ms for an unknown
resolution. Bus overhead adds to the nominal budget.
\begin{lstlisting}
float temperature_c = 0.0f;
if (address_ok) {
    result = DS18B20_ReadTemperatureBlocking(
        &bus, rom, 1000U, &temperature_c);
}
bool temperature_valid = (result == DS18B20_STATUS_OK);
\end{lstlisting}
This continues the preceding example. Use the output only when valid; initial
zero and old values are not successful measurements. This function can block
for roughly a conversion interval. For concurrent main-loop work use the array
interface or the split sequence below.

\subsection{Split conversion from reading}
\texttt{DS18B20\_StartConversion} sends Convert T to a ROM; null broadcasts to the
bus. OK means command sent, not measurement completed.
\texttt{DS18B20\_IsConversionReady} checks the conversion busy bit.
\texttt{DS18B20\_ReadTemperature} reads a completed conversion and performs CRC
and range checks; it does not start a new measurement. Its output changes only
on success.

On STM32 use the HAL tick; elsewhere supply a monotonic millisecond clock.
Start once and retain the state variables between main-loop iterations:
\begin{lstlisting}
bool waiting = false;
uint32_t deadline = 0;
temperature_valid = false;
result = DS18B20_StartConversion(&bus, rom);
if (result == DS18B20_STATUS_OK) {
    deadline = HAL_GetTick() + 760U; /* After command. */
    waiting = true;
}
\end{lstlisting}
Run this fragment repeatedly in the main loop:
\begin{lstlisting}
if (waiting &&
    (int32_t)(HAL_GetTick() - deadline) >= 0) {
    waiting = false;
    result = DS18B20_IsConversionReady(&bus)
        ? DS18B20_ReadTemperature(&bus, rom, &temperature_c)
        : DS18B20_STATUS_TIMEOUT;
    temperature_valid = (result == DS18B20_STATUS_OK);
}
/* Other main-loop work continues while waiting. */
\end{lstlisting}
Do not reset/search or send another bus command between conversion and polling.
The array scheduler handles this sequence for you and is usually simpler.

\subsection{Prepare buffers: ArrayInit}
Inputs are manager, initialized bus, mutable ROM array, separate discovery
buffer, temperature array, status array and count. All buffers have that capacity
and remain alive while sampling. Invalid pointers/count, duplicate ROMs and
invalid nonzero ROMs cause false. Zero ROMs request automatic discovery.
\begin{lstlisting}
#define N 4U /* Use 1U for one sensor. */
static DS18B20_Array sensors;
static DS18B20_Rom addresses[N] = {{0}};
static DS18B20_Rom discovered[N];
static float temperatures[N];
static DS18B20_Status status[N];
bool initialized = DS18B20_ArrayInit(
    &sensors, &bus, addresses, discovered,
    temperatures, status, N);
if (initialized) sensors.sample_interval_ms = 250U;
\end{lstlisting}
Run once, not every loop iteration. The interval is the idle gap after a cycle,
not conversion duration. For known ROMs fill the address array before initialization
in wire order: family, six serial bytes, CRC. Discovery/replacement can change
this array, so do not declare it const. Static buffers avoid large stack use.

\subsection{Service the array: ArrayProcess}
Inputs are manager and current millisecond tick. Call frequently. True means a
complete cycle ended, even if some slots failed; false means no completed cycle
yet. Check each slot's status separately. Index \texttt{i} always corresponds
to the same address-array index.
\begin{lstlisting}
static float latest_c[N];
static bool valid[N];
if (initialized &&
    DS18B20_ArrayProcess(&sensors, HAL_GetTick())) {
    for (size_t i = 0; i < N; ++i) {
        valid[i] = (status[i] == DS18B20_STATUS_OK);
        if (valid[i]) latest_c[i] = temperatures[i];
    }
}
\end{lstlisting}
This optional snapshot preserves the last numbers but marks failed slots invalid.
Use a value for display/control only when its valid flag is true. The library's
own outputs become NAN during conversion or errors.

\subsection{Replace a chosen slot: ArrayReplace}
Inputs are manager, zero-based destination index and new valid ROM. True updates
the mapping and invalidates the old temperature. False indicates a bad ROM/index,
a ROM already used by another slot or a non-idle array. Obtain the new ROM from
a complete scan or a separate single-sensor setup, not Read ROM on a multidrop bus.
\begin{lstlisting}
/* replacement_rom was obtained and validated earlier. */
if (sensors.phase == 0U) {
    bool replaced = DS18B20_ArrayReplace(
        &sensors, 1U, replacement_rom); /* Second slot. */
    /* Check replaced, then wait for a valid new sample. */
}
\end{lstlisting}
Use at least two slots for this example and call from the same task as processing.
ROM changes do not survive reset unless the application saves them. The following
section explains when replacement can be automatic.

\subsection{Configure and verify resolution}
\texttt{DS18B20\_WriteScratchpad} takes bus, ROM, high alarm, low alarm and resolution.
The example sets 10-bit resolution and alarms at 30 and 10 C. Settings are volatile;
no EEPROM copy is performed. A successful command alone is not an addressed sensor
acknowledgment, so verify using \texttt{DS18B20\_ReadScratchpad} and
\texttt{DS18B20\_GetResolution}.
\begin{lstlisting}
result = DS18B20_WriteScratchpad(
    &bus, rom, 30, 10, DS18B20_RESOLUTION_10_BIT);
uint8_t pad[DS18B20_SCRATCHPAD_SIZE];
if (result == DS18B20_STATUS_OK)
    result = DS18B20_ReadScratchpad(&bus, rom, pad);
bool configured = (result == DS18B20_STATUS_OK) &&
    DS18B20_GetResolution(pad) == DS18B20_RESOLUTION_10_BIT &&
    pad[2] == 30U && pad[3] == 10U;
\end{lstlisting}
Configure before sampling or while idle. The array still uses its conservative
12-bit wait. \texttt{DS18B20\_GetConversionTimeMs} returns 94, 188, 375 or 750 ms
for 9--12 bits; add timing margin when designing your own scheduler.

\subsection{Raw helpers and errors}
\texttt{DS18B20\_ReadScratchpad} reads and validates nine bytes.

After a successful read, \texttt{DS18B20\_DecodeRaw} extracts the signed raw
temperature and masks undefined bits.

\texttt{DS18B20\_RawToCelsius} converts the raw value to degrees.

For normal measurements, use \texttt{ReadTemperature}. It also rejects values
outside the sensor's valid range.

OK permits use; BUSY means wait; NO\_DEVICE means no reset presence; CRC\_ERROR
means unusable data; TIMEOUT means conversion was not ready; OUT\_OF\_RANGE means
an impossible value; INVALID\_ARGUMENT means correct the call or ROM. A missing
addressed sensor can produce CRC\_ERROR if another sensor answers reset. Always
check status, not just the numeric value. The OneWire reset/bit/byte functions
are protocol primitives usually called by the sensor layer; their implementation
remains in the source files rather than being copied into this report.
